% Fianza whitepaper
% Compile:  pdflatex whitepaper.tex   (run twice for refs/TOC)
% Overleaf: upload as-is, no extra files needed. Standard packages only.

\documentclass[11pt,a4paper]{article}

\usepackage[T1]{fontenc}
\usepackage[utf8]{inputenc}
\usepackage{lmodern}
\usepackage[margin=1.15in]{geometry}
\usepackage{amsmath,amssymb}
\usepackage{booktabs}
\usepackage{array}
\usepackage{microtype}
\usepackage{titlesec}
\usepackage{enumitem}
\usepackage[hidelinks]{hyperref}

% --- look & feel -------------------------------------------------------------
\linespread{1.06}
\setlength{\parskip}{0.45em}
\setlength{\parindent}{0pt}

\titleformat{\section}{\normalfont\large\bfseries}{\thesection.}{0.6em}{}
\titleformat{\subsection}{\normalfont\normalsize\bfseries}{\thesubsection}{0.6em}{}
\titlespacing*{\section}{0pt}{1.6em}{0.6em}
\titlespacing*{\subsection}{0pt}{1.1em}{0.4em}

% clamp operator used throughout the independence model
\DeclareMathOperator{\clampop}{clamp}
\newcommand{\clamp}[1]{\clampop_{[0,1]}\!\left(#1\right)}

\title{\vspace{-1.2cm}\textbf{Fianza: Uncollateralized Credit for Autonomous Agents}\\[0.35em]
\large A revenue-underwritten lending protocol on Stellar}

\author{Divyanshh Kalra \and Kundan Kumar}
\date{Version 1.0 \quad\textbullet\quad August 2026 \quad\textbullet\quad \texttt{docs.fianza.space}}

\begin{document}
\maketitle

% --- abstract ----------------------------------------------------------------
\begin{abstract}
\noindent
An autonomous software agent can now earn money. It sells inference, research and
data over machine-payable APIs and is paid in stablecoins, on-chain, without a
human in the loop. It still cannot borrow. No bank underwrites a process, and
collateralized lending requires it to already hold the capital it is trying to
raise. The result is a class of economically productive actors that can generate
income but cannot finance the inputs that income depends on.

We propose a lending protocol that underwrites an agent's \emph{verifiable
trailing revenue} rather than its assets. The central difficulty is not reading
revenue---inflows are public---but establishing that revenue represents genuine
third-party economic activity rather than an operator paying themselves from
wallets they control. We show that this cannot be \emph{decided}, only
\emph{priced}, and we give a continuous discount model that converts raw inflows
into effective independent revenue by weighting each payer on account age,
external counterparty diversity, funding-graph independence and net-flow
reciprocity, then penalising concentration. Credit is sized as a tier multiple of
that figure, and further restricted by a repayment-history ramp beginning at
15\% of the sized limit. Each borrower's lenders occupy an isolated vault, so a
default is contained to the parties who chose that borrower. Defaults are
triggered permissionlessly, absorbed first by a reserve funded from interest, and
otherwise socialised pro-rata through share price. The system is implemented in
three Soroban contracts settling in USDC and is live on Stellar mainnet.
\end{abstract}

\vspace{0.6em}
\hrule
\vspace{0.4em}

% --- 1 -----------------------------------------------------------------------
\section{Introduction}

Machine-payable APIs made agents into economic participants. An agent can hold a
keypair, receive stablecoin payments for work performed, and pay for its own
inputs. Payment is solved. Credit is not.

Consider an agent that reliably earns \$5 for a research task requiring a \$2
data purchase. It is profitable. If its balance is below \$2 at the moment work
arrives, it cannot begin, and the profit is lost. Traditional finance cannot
serve it: underwriting assumes a legal person, a jurisdiction, and recourse.
Collateralized DeFi cannot serve it either---a protocol requiring \$150 of
collateral to lend \$100 is useless to a borrower whose defining characteristic
is having no capital and a real income stream.

The gap is not liquidity. It is underwriting. What is missing is a way to convert
observable income into a credit decision a contract can enforce.

Three properties make this newly tractable. Agent revenue arrives on a public
ledger, so it is observable without self-reporting. Agents are programs, so
borrowing and repayment can be automated to a degree no human borrower matches.
And settlement on a low-fee, fast-finality chain makes loans of a few dollars
economically coherent, where elsewhere gas alone would exceed the interest.

Section~\ref{sec:primitive} states the underwriting primitive.
Section~\ref{sec:independence} addresses the adversarial core: distinguishing
independent revenue from manufactured revenue. Sections~\ref{sec:scoring}
through~\ref{sec:default} give the credit, custody and default mechanics.
Section~\ref{sec:security} analyses economic security, and
Section~\ref{sec:limits} states the limits plainly.

% --- 2 -----------------------------------------------------------------------
\section{Prior approaches and why they do not apply}

\textbf{Collateralized lending.} Over-collateralized protocols avoid
underwriting entirely: the loan is safe because the collateral exceeds it. This
is robust and useless here, since it presupposes the capital the borrower lacks.

\textbf{Credit scoring from wallet history.} Several systems score addresses on
transaction history, protocol interaction or asset holdings. These measure
\emph{activity}, not \emph{income}, and activity is cheap to manufacture. A score
that can be farmed is not an underwriting signal.

\textbf{Off-chain underwriting with on-chain settlement.} Real-world credit
protocols underwrite legal entities off-chain and settle on-chain. This works,
and does not generalise to agents: there is no entity to underwrite, no
jurisdiction, and no recourse.

\textbf{Reputation and staking.} Requiring a bond makes misbehaviour costly but
again requires capital up front, reproducing the original problem.

Our approach differs in what it treats as the primitive: not the borrower's
assets, identity or reputation, but the \emph{cash flows arriving at its
address}, and the degree to which those flows can be shown to originate from
parties the borrower does not control.

% --- 3 -----------------------------------------------------------------------
\section{Revenue as an underwriting primitive}
\label{sec:primitive}

Let an agent hold address $a$. Over a trailing window we observe the set of
stablecoin transfers into $a$. Each transfer has a payer, an amount and a
timestamp, all publicly verifiable. Define raw revenue as
\begin{equation}
R_{\mathrm{raw}} \;=\; \sum_{i \in P} r_i
\end{equation}
where $P$ is the set of distinct payers and $r_i$ the total received from payer
$i$.

$R_{\mathrm{raw}}$ is worthless as an underwriting input. It is trivially
inflated: an operator creates $n$ addresses, funds them, and cycles capital
through them. The entire problem is converting $R_{\mathrm{raw}}$ into a figure
reflecting only revenue the operator could not have manufactured cheaply.

We additionally admit off-chain revenue proven by zero-knowledge TLS
attestation---an agent proving a payment-processor balance without revealing
credentials---verified by an on-chain verifier contract. Proven off-chain revenue
is weighted at $1.5\times$ relative to raw on-chain revenue, since it carries an
external attestation on-chain revenue lacks.

% --- 4 -----------------------------------------------------------------------
\section{The independence problem}
\label{sec:independence}

\subsection{Statement}

An underwriter must answer: \emph{is this revenue genuine independent economic
activity, or the operator paying itself?}

This admits no proof. Nothing distinguishes, at the ledger level, a payment from
a real customer from a payment by the operator's second wallet. Both are valid
transfers of real value. The distinguishing fact---common control---is not
observable.

We therefore restate the goal. We do not attempt to detect fake revenue. We
attempt to make manufacturing \emph{counted} revenue more expensive than the
credit it unlocks. The protocol is solvent by construction if
\begin{equation}
\mathrm{cost\_to\_fake}(R) \;>\; \mathrm{value\_unlocked}(R)
\label{eq:invariant}
\end{equation}
Since credit is at most $3\times$ counted revenue and the ramp restricts a cold
borrower to 15\% of that, an attacker must cycle real capital repeatedly through
apparently unrelated wallets to extract a small multiple. Every mechanism below
exists to raise the left side of~\eqref{eq:invariant}.

\subsection{Threat model}

We assume an adversary who can create unlimited addresses at negligible cost,
move their own capital freely between them, wait to age accounts, script
plausible activity, and obtain genuine external payment-processor accounts. We
assume they \emph{cannot} cheaply obtain other people's independent economic
activity. That single asymmetry is the foundation of the model.

\begin{center}
\begin{tabular}{@{}llp{7.6cm}@{}}
\toprule
\textbf{ID} & \textbf{Attack} & \textbf{Mechanism} \\
\midrule
A1 & Self-payment      & Agent pays itself from controlled wallets \\
A2 & Fresh-wallet farm & $n$ new addresses with no other history \\
A3 & Circular funding  & Agent funds wallets, they pay it back \\
A4 & Concentration     & Two or three payers supply nearly all revenue \\
A5 & Synthetic cadence & Scripted, regular payments imitating volume \\
A6 & Off-chain self-charge & Operator charges their own cards \\
A7 & Collusion ring    & Distinct real operators cross-pay to inflate each other \\
\bottomrule
\end{tabular}
\end{center}

\subsection{Per-payer weighting}

Each payer $i$ receives a weight $w_i \in [0,1]$, the product of four independent
factors. A product, not a sum: any single factor going to zero eliminates that
payer's contribution entirely, which is the correct behaviour when a payer is
provably circular.
\begin{equation}
w_i \;=\; \alpha_i \cdot \delta_i \cdot \phi_i \cdot \rho_i
\end{equation}

\textbf{Age.} $\alpha_i = \clamp{\dfrac{\mathrm{age}_i}{A}}$, where
$\mathrm{age}_i$ is the payer's account age and $A$ a maturity threshold.
Defeats A2: a wallet created to pay the agent contributes nothing.

\textbf{External diversity.} $\delta_i = \clamp{\dfrac{d_i}{D}}$, where $d_i$ is
the number of distinct counterparties the payer transacts with \emph{excluding
the agent and the agent's other payers}, and $D$ a diversity threshold. Defeats
A1 and weakens A7: a wallet that only ever deals with this agent is a puppet; one
that deals only with the agent and its co-payers is a ring member. Both show zero
external out-degree.

\textbf{Funding independence.} $\phi_i \in \{0,1\}$, zero if the payer's funds
trace back to the agent within $K$ hops of the transfer graph, one otherwise.
Defeats A3 directly. We perform a breadth-first traversal of funding sources to
depth $K = 3$.

\textbf{Reciprocity.} $\rho_i = \clamp{1 - \dfrac{p_i}{r_i}}$, where $p_i$ is
what the agent paid \emph{back} to payer $i$. Mutual cross-payment nets toward
zero, attacking A7: a ring whose members pay each other in both directions
produces little counted revenue.

\subsection{Concentration}

Even with genuine, independent payers, revenue concentrated in one or two
counterparties is fragile and easier to fabricate. Two mechanisms apply. First,
each payer's contribution is capped at a maximum share $\theta$ of raw revenue
before weighting:
\begin{equation}
c_i \;=\; \min\!\left(r_i,\; \theta \cdot R_{\mathrm{raw}}\right)
\end{equation}
Second, we compute a normalised Herfindahl--Hirschman index over capped shares
$s_i = c_i / \sum_j c_j$:
\begin{equation}
H \;=\; \sum_i s_i^{\,2}
\end{equation}
and derive a concentration factor that is $1$ for well-diversified revenue and
decays to $0$ as revenue concentrates:
\begin{equation}
\gamma \;=\; 1 - \clamp{\frac{H - H_0}{1 - H_0}}
\end{equation}
where $H_0$ is a tolerance floor below which no penalty applies. This addresses
A4.

\subsection{Temporal organicity}

Genuine demand is irregular. Scripted payments are not. We compute a factor
$\omega \in [0,1]$ from the dispersion of inter-payment intervals, penalising
distributions that are implausibly uniform. This raises the cost of A5 without
claiming to defeat a patient adversary who deliberately randomises.

\subsection{Aggregation}

Effective independent revenue is
\begin{equation}
R_{\mathrm{eff}} \;=\; \gamma \cdot \omega \cdot \sum_{i \in P} w_i \, c_i
\label{eq:reff}
\end{equation}

$R_{\mathrm{eff}}$, never $R_{\mathrm{raw}}$, is what enters scoring. The ratio
$R_{\mathrm{eff}} / R_{\mathrm{raw}}$ is published as an independence score,
making the discount auditable rather than opaque.

The output is a discount, not a verdict. A payer is not accepted or rejected; it
is worth some fraction of its face value. This matters: a binary classifier has a
decision boundary an attacker can probe and cross. A continuous discount means
that as an attacker's construction approaches indistinguishability from real
revenue, its cost approaches that of real revenue.

% --- 5 -----------------------------------------------------------------------
\section{Scoring and credit sizing}
\label{sec:scoring}

$R_{\mathrm{eff}}$ and the count of independent counterparties map to a composite
score $S \in [0, 850]$ via revenue bands plus a diversity bonus of 5 points per
independent payer up to 10 payers. On-chain revenue counts toward the score only
if the agent clears a minimum counterparty threshold; below it, and absent proven
off-chain revenue, the score is capped below the lending threshold.

The score bands into a tier, and the tier fixes both the credit multiple and the
base interest rate.

\begin{center}
\begin{tabular}{@{}lccc@{}}
\toprule
\textbf{Tier} & \textbf{Score} & \textbf{Limit multiple} & \textbf{Base APR} \\
\midrule
A       & $\geq 750$ & $3.0\times$ & 6.00\% \\
B       & $\geq 650$ & $2.0\times$ & 8.50\% \\
C       & $\geq 550$ & $1.0\times$ & 12.00\% \\
Unrated & $< 550$    & $0$         & --- \\
\bottomrule
\end{tabular}
\end{center}

The sized limit is
\begin{equation}
L_{\mathrm{base}} \;=\; R_{\mathrm{eff}} \times m(\mathrm{tier})
\end{equation}
An agent with no revenue scores Unrated and receives a zero limit. This is the
correct output, not a failure state.

The multiple is deliberately small. At $3\times$ for the best tier, an attacker
must manufacture \$1 of \emph{counted} revenue to unlock at most \$3 of exposure,
before the ramp reduces it further.

% --- 6 -----------------------------------------------------------------------
\section{The credit ramp}
\label{sec:ramp}

A score derived from revenue says what an agent \emph{could} support. It says
nothing about whether it repays. We therefore separate the sized limit from the
enforceable limit.

Let $n$ be on-time repayments and $k$ missed. The ramp factor is
\begin{equation}
f \;=\; \clamp{f_0 + n\,\Delta^{+} - k\,\Delta^{-}}
\end{equation}
with $f_0 = 0.15$, $\Delta^{+} = 0.15$ and $\Delta^{-} = 0.30$. The enforced
limit is
\begin{equation}
L \;=\; L_{\mathrm{base}} \cdot f
\end{equation}

Three consequences. A cold agent, however well-scored, draws at most 15\% of its
sized limit, bounding what a freshly-fabricated identity can extract. Trust is
earned linearly and lost at twice the rate. And a degrading agent manifests as a
ramp that stops advancing---an observable leading indicator preceding default
rather than following it.

This is the mechanism that makes the system robust to the volatility of its
borrowers. An agent's competence may change without warning: its underlying model
may be swapped, its market may move. We do not attempt to predict this. We bound
exposure at any moment and let history move the bound.

% --- 7 -----------------------------------------------------------------------
\section{Isolated vaults}

Each agent has its own vault. A lender does not deposit into a protocol-wide
pool; they deposit against a specific borrower, receiving shares in that vault
alone. For a vault with total assets $A_v$ and total shares $T$, a deposit of $x$
mints
\begin{equation}
\Delta T \;=\;
\begin{cases}
x, & T = 0 \\[0.3em]
x \cdot T / A_v, & \text{otherwise}
\end{cases}
\end{equation}
Share price $A_v/T$ therefore rises with accrued interest and falls with realised
loss, and every lender's position tracks the performance of the borrower they
selected.

Isolation is a deliberate rejection of capital efficiency in favour of
containment. A shared pool is more efficient and spreads yield; it also spreads
\emph{loss}, and for a borrower class where high default rates are expected
rather than exceptional, that is the wrong trade at this stage. A pooled layer
can be added above isolated vaults once loss rates are empirically characterised;
it cannot be safely removed later.

% --- 8 -----------------------------------------------------------------------
\section{Interest, reserve and yield}

A draw starts a term. Interest accrues linearly on outstanding principal at
\begin{equation}
\mathrm{APR} \;=\; \mathrm{APR}_{\mathrm{base}}(\mathrm{tier}) + \pi \cdot u
\end{equation}
where $u = \mathrm{principal} / \mathrm{assets}$ is vault utilisation and $\pi$
the maximum utilisation premium (10 percentage points at full utilisation).
Scarce liquidity is priced upward, which both compensates lenders and dampens
demand before a vault is drained.

Repayment applies to interest before principal. Interest is split: a fraction
$\sigma = 20\%$ accrues to a first-loss \emph{reserve}, and the remainder becomes
claimable lender \emph{yield}:
\begin{equation}
\mathrm{reserve} \mathrel{+}= \sigma I,
\qquad
\mathrm{yield} \mathrel{+}= (1-\sigma) I
\end{equation}
The reserve is the mechanism by which performing loans pay for failing ones. It
is funded by borrowers, not lenders, and accumulates in the specific vault whose
borrower generated it.

% --- 9 -----------------------------------------------------------------------
\section{Default}
\label{sec:default}

If principal remains outstanding past the due date, \textbf{any} party may call
\texttt{mark\_default}. Permissionlessness matters: the protocol operator is not
a required participant in loss realisation, and cannot delay recognition of a bad
loan.

On default, with outstanding principal $D$ and reserve $V$:
\begin{equation}
V_{\mathrm{used}} = \min(D, V),
\qquad
\Lambda = D - V_{\mathrm{used}}
\end{equation}
The reserve is consumed first and returns to withdrawable liquidity. The
remainder $\Lambda$ is written off as realised loss. Because total shares are
unchanged while assets fall by $\Lambda$, the loss is socialised pro-rata through
share price automatically---no iteration over lenders, no privileged distribution
step. Accrued interest is written off, the agent is frozen from further
borrowing, and the miss enters its repayment record, collapsing the ramp by
$\Delta^{-}$ and reducing its score at the next underwriting pass.

The loss is bounded by the vault. Other agents' vaults are untouched, by
construction rather than by policy.

% --- 10 ----------------------------------------------------------------------
\section{Incentives}

\textbf{Borrowers} repay because the ramp is the only path to a larger line, and
because default freezes borrowing permanently and is recorded on-chain against
the address. For an agent whose business depends on working capital, the option
value of future credit exceeds a single stolen draw once the ramp has advanced.
The system's weakest moment is a cold agent's first draw, which is precisely what
$f_0 = 0.15$ bounds.

\textbf{Lenders} earn tier-priced interest plus a utilisation premium, choose
their borrower explicitly, and are protected by a reserve their borrower funded.
They are not asked to trust agents in general---only the one they selected.

\textbf{Attackers} face inequality~\eqref{eq:invariant}. Each mechanism raises
$\mathrm{cost\_to\_fake}$: aging wallets costs time, building external
counterparty diversity costs genuine relationships, avoiding the funding graph
costs an unlinked capital source, and defeating concentration costs breadth. None
is individually insurmountable; the model's claim is that jointly they cost more
than $3f$ per unit of counted revenue.

% --- 11 ----------------------------------------------------------------------
\section{Security analysis}
\label{sec:security}

Consider an adversary targeting the simplest complete attack: fabricate revenue,
draw the maximum, and abscond.

To be counted at all, each fabricated payer must be aged past $A$, transact with
$\lceil D \rceil$ external counterparties that are not the agent's other payers,
and be funded from a source not traceable to the agent within $K$ hops. To avoid
the concentration penalty the adversary needs several such payers, since a single
payer is capped at $\theta$ of raw revenue and drives $H \to 1$, sending
$\gamma \to 0$.

Suppose they assemble this and achieve $R_{\mathrm{eff}} = R$. Tier A requires a
high score, so realistically they land in tier B or C, giving
$L_{\mathrm{base}} \leq 2R$. Cold, the ramp restricts them to
$0.15 \cdot L_{\mathrm{base}} \leq 0.3R$. To do better they must \emph{repay}---
returning capital and paying interest---several times, each cycle advancing the
ramp by 15\%. Reaching the full limit takes six on-time repayments.

The adversary's maximum theft is therefore bounded by roughly $2R$ after
sustained good behaviour, having cycled real capital through independent-looking
wallets throughout, and having paid interest that funded the reserve now
absorbing part of their default. Whether this inequality holds in practice is an
empirical question only mainnet loss data can settle; we state the target rather
than claim it is met.

\textbf{The open weakness.} A sophisticated, non-reciprocal collusion ring---
several genuinely distinct operators, with real external activity, each paying
one direction only---is not distinguishable from real customers at the graph
level. Reciprocity nets out mutual rings; external diversity catches closed ones.
A ring that is neither mutual nor closed defeats the current model. We believe
this requires either staking, which reintroduces a capital requirement, or global
community detection over a longitudinal graph. We consider this unsolved and
state it rather than obscure it.

% --- 12 ----------------------------------------------------------------------
\section{Implementation}

The protocol runs as three Soroban contracts on Stellar, settling in USDC via its
SEP-41 contract interface.

\begin{description}[leftmargin=1.3em,style=unboxed,itemsep=0.3em]
\item[\texttt{score\_registry}] stores each agent's published score, tier,
revenue basis and repayment record. Publication is signer-authorised; reads are
open. It is the single source of truth other contracts consult.

\item[\texttt{credit\_line}] is a read-only view deriving terms---tier, ramped
limit, APR---from a published score, using the same shared policy library the
vault enforces. Quoted terms and enforced terms cannot diverge, because they are
the same code.

\item[\texttt{lending\_vault}] holds custody. Deposits, share accounting,
borrows, the repayment waterfall, reserve, yield, permissionless default and loss
socialisation all live here. It independently re-derives the limit from the
registry rather than trusting a caller.
\end{description}

A shared \texttt{revenue\_math} library holds all policy---banding, multiples,
APR, ramp, interest split, utilisation---and is compiled into every contract that
needs it. This is the invariant keeping the system coherent: there is one
definition of each credit rule.

Underwriting itself is off-chain: indexing transfer events, traversing the
funding graph, computing $R_{\mathrm{eff}}$, and producing a signed score. The
computation is heavy and iterative, and is not a good fit for on-chain execution.
Its \emph{output} is on-chain and its \emph{inputs} are public, so any third
party can recompute a published score from ledger data and detect a dishonest
underwriter. Trust is minimised by verifiability rather than by execution
location.

Stellar is load-bearing here, not incidental. Loans in this system are often
under ten dollars. On a chain where a settlement costs more than the interest on
such a loan, the product does not exist. Sub-cent fees and roughly five-second
finality are what make per-draw on-chain settlement coherent at this size.

% --- 13 ----------------------------------------------------------------------
\section{Limitations}
\label{sec:limits}

We state these plainly, since a credit protocol that hides its failure modes is
not one worth lending through.

\textbf{A single-signer trust root.} Scores are published by one authorised
signer today. A compromised signer could publish fraudulent scores. Scores are
recomputable from public data, so fraud is detectable, but detection is not
prevention. Multi-signature attestation is the necessary next step.

\textbf{Collusion rings are unsolved}, as described in
Section~\ref{sec:security}.

\textbf{Off-chain revenue is partially trusted.} A zkTLS proof establishes that a
balance exists, not that it represents independent income. An operator charging
their own cards produces a valid proof (A6). We weight proven off-chain revenue
above raw on-chain revenue, which is defensible for authenticity and generous for
independence.

\textbf{Parameters are calibrated on thin data.} The thresholds $A$, $D$,
$\theta$, $H_0$ and the ramp constants are chosen from reasoning and adversarial
testing, not from observed default behaviour at scale, which does not yet exist.

\textbf{Agent behaviour is not modelled.} We deliberately do not attempt to
predict whether an agent will keep performing. Our protection is exposure
bounding---short terms, small multiples, a ramp---not forecasting. If a genuinely
earning agent degrades, the loss is real and lands on that vault's lenders.

% --- 14 ----------------------------------------------------------------------
\section{Conclusion}

We have described a lending protocol that underwrites autonomous agents on
verifiable revenue rather than collateral. The contribution is not the vault, the
interest curve or the default path---those are standard. It is the treatment of
counterparty independence as an economic quantity rather than a classification
problem: revenue is discounted continuously by how expensive it would have been
to manufacture, and credit is sized against what survives that discount.

We claim no proof of independence. We claim a construction in which faking
counted revenue costs more than the credit it unlocks, bounded further by a
repayment ramp and contained by per-borrower isolation. Whether the inequality
holds under a real adversary with real money is the question the next phase must
answer, and it can only be answered on mainnet.

The system is live. The contracts are deployed on Stellar mainnet and testnet,
the loop has executed end to end in real USDC, and a default has been triggered
and absorbed on-chain exactly as specified here.

% --- refs --------------------------------------------------------------------
\begin{thebibliography}{9}
\setlength{\itemsep}{0.15em}
\bibitem{btc} S. Nakamoto. \emph{Bitcoin: A Peer-to-Peer Electronic Cash System.} 2008.
\bibitem{x402} Coinbase et al. \emph{x402: An HTTP-native payment protocol.} 2025.
\bibitem{soroban} Stellar Development Foundation. \emph{Soroban smart contract platform.}
\bibitem{sep41} Stellar Development Foundation. \emph{SEP-41: Token Interface.}
\bibitem{reclaim} Reclaim Protocol. \emph{zkTLS attestations for web data.}
\bibitem{hhi} A. Hirschman. \emph{The Paternity of an Index.} American Economic Review, 1964.
\bibitem{sybil} Fianza. \emph{Sybil and counterparty-independence model.} \texttt{docs.fianza.space/sybil-model}
\bibitem{engine} Fianza. \emph{Credit engine.} \texttt{docs.fianza.space/credit-engine}
\end{thebibliography}

\appendix

% --- A -----------------------------------------------------------------------
\section{Parameters}

\begin{center}
\begin{tabular}{@{}llr@{}}
\toprule
\textbf{Symbol} & \textbf{Meaning} & \textbf{Value} \\
\midrule
$K$          & Funding-graph traversal depth        & 3 hops \\
$\theta$     & Max single-payer share before weighting & 0.40 \\
$H_0$        & HHI tolerance floor                  & 0.15 \\
$f_0$        & Ramp start                           & 15\% \\
$\Delta^{+}$ & Ramp gain per on-time repayment      & $+15\%$ \\
$\Delta^{-}$ & Ramp penalty per miss                & $-30\%$ \\
$\sigma$     & Interest share to reserve            & 20\% \\
$\pi$        & Max utilisation premium              & $+10.00\%$ \\
---          & Tier A / B / C multiple              & $3.0\times$ / $2.0\times$ / $1.0\times$ \\
---          & Tier A / B / C base APR              & 6.00\% / 8.50\% / 12.00\% \\
---          & Minimum independent counterparties   & 3 \\
---          & Off-chain (zkTLS) revenue weight     & $1.5\times$ \\
\bottomrule
\end{tabular}
\end{center}

Age and diversity thresholds $A$ and $D$ are network-calibrated; the values used
on testnet are deliberately lower than mainnet targets because the chain and the
agents on it are young, and are documented in the Sybil model reference.

% --- B -----------------------------------------------------------------------
\section{Worked example}

An agent has received \$120 from four payers: \$60, \$30, \$20, \$10.

\textbf{Capping.} $R_{\mathrm{raw}} = 120$, so $\theta R_{\mathrm{raw}} = 48$.
The largest payer is capped: $c = (48, 30, 20, 10)$, $\sum c = 108$.

\textbf{Weighting.} Payer 1 is mature and externally diverse, $w_1 = 1$. Payer 2
was funded by the agent two hops back, so $\phi_2 = 0$ and $w_2 = 0$. Payer 3 is
fresh at half the age threshold with adequate diversity, $w_3 = 0.5$. Payer 4 is
mature but the agent paid \$4 back, so $\rho_4 = 1 - 4/10 = 0.6$ and $w_4 = 0.6$.
\[
\sum_i w_i c_i \;=\; 48 + 0 + 10 + 6 \;=\; 64
\]

\textbf{Concentration.} Shares over capped total:
$(0.4444,\, 0.2778,\, 0.1852,\, 0.0926)$, giving $H = 0.3176$. With $H_0 = 0.15$:
\[
\gamma \;=\; 1 - \frac{0.3176 - 0.15}{0.85} \;=\; 0.8029
\]

\textbf{Aggregation.} Assuming organicity $\omega = 1$:
\[
R_{\mathrm{eff}} \;=\; 0.8029 \times 64 \;\approx\; 51.38
\]
The agent presents \$120 of inflows and is underwritten on \$51.38---an
independence score of $0.428$.

\textbf{Sizing.} Suppose this scores into tier C:
$L_{\mathrm{base}} = 51.38 \times 1.0 = 51.38$. With no repayment history,
$f = 0.15$, so the enforceable limit is
\[
L \;=\; 51.38 \times 0.15 \;\approx\; \$7.71
\]

From \$120 of raw inflows, a cold agent may draw under eight dollars. Six on-time
repayments bring the ramp to 100\%, at which point the full \$51.38 becomes
available. That gap between apparent and extractable credit is the protocol's
primary defence.

\end{document}
